随着高校毕业论文、课程论文及各类学术论文写作任务的增多，论文排版的规范性与编辑效率逐渐成为学生和科研人员关注的重要问题。传统文字处理工具虽然操作直观，但在复杂公式、图表编号、交叉引用、参考文献管理和论文模板适配等场景下，仍需大量人工调整；LaTeX 虽然具有较强的学术排版能力，但其命令语法、环境结构、编译配置和错误日志分析对初学者并不友好。针对普通用户使用 LaTeX 时存在的代码记忆成本高、模板使用复杂、编译错误难以定位等问题，本文设计并实现了一套基于 JavaWeb 和 LaTeX 的学术论文智能排版系统，为论文写作与排版提供在线化、组件化和智能化支持。

系统采用前后端分离架构，前端基于 Vue 3 和 Vite 实现页面展示与交互功能，后端基于 Spring Boot、MyBatis 和 MySQL 完成业务逻辑处理与数据持久化。围绕学术论文写作流程，系统实现了用户注册登录、论文项目管理、LaTeX 在线编辑、模板管理、论文编译、PDF 预览、文件导出、AI 对话辅助和编译错误分析等功能。用户可直接在浏览器中创建论文项目，选择或导入论文模板，在线编辑论文内容，并通过不同编译引擎生成 PDF 文档。

本文系统的核心设计在于引入面向论文结构的组件库拖拽式插入机制。系统将章节、图片、表格、公式、列表、符号等常用 LaTeX 论文结构封装为可视化组件，用户无需专门记忆完整的 LaTeX 命令和环境写法，只需从组件库中选择所需组件并拖入编辑区域，系统即可自动生成相应的 LaTeX 代码框架。该方式降低了 LaTeX 语法使用门槛，减少了命令拼写错误、括号缺失和环境未闭合等常见问题，提高了论文编辑效率。

在具体实现过程中，系统还针对 LaTeX 模板依赖文件较多、编译错误信息不易理解等问题进行了优化。系统支持 zip 模板包导入，并在模板解析和编译过程中尽量保留原有目录结构，减少因 .cls、.sty、.bib 等依赖文件缺失导致的编译失败。同时，系统引入 AI 辅助功能，能够结合编译错误日志给出原因分析和修改建议，帮助用户快速定位问题。测试结果表明，系统能够完成论文项目管理、组件拖拽插入、模板载入、在线编辑、论文编译、PDF 预览、文件导出和 AI 错误分析等主要任务，整体运行较为稳定，基本满足学术论文在线排版需求，具有一定的实际应用价值。